\subsubsection*{The span of the original theta rows}
\label{subsec:theta-row-rank}

Integrality after descent bounds the coefficients of the theta rows. For their determinant to detect the original row space, these rows must also span its good type at the base tuple with the prescribed unipotent character.

\paragraph{Stable range and the lattice operator.}
Write \(n\) for the source rank, \(r\) for the target
rank and \(D_E=[E:\mathbb Q]r\) for the dimension of a
target lattice over \(\mathbb Q\). The source dimension
condition is \(n\ge 4D_E+4\), with \(n\) even as in the
row construction. In the finite Weil model
for the prescribed level, let \(\Phi(L,A)\) denote the
normalized test vector for a self-dual target lattice
\(L\) and a source lattice class \(A\). The test vectors
are summed over the finite source classes with their
stabilizer weights. Their pairings define the matrix
\[
      C(L,L')=\sum_A w(A)
                \langle\Phi(L,A),\Phi(L',A)\rangle .
\]
With the same weights on both entries, its diagonal
is a positive scalar. If \(m=[L':L\cap L']\), the
ratio of the off-diagonal bound to the diagonal is
\(m^{-n/2}\). We give the common-label calculation
that makes this bound independent of the size of the level.

\begin{lemma}[Invertibility of the overlap matrix in the stable range]\label{lem:theta-overlap-invertible}
At the specified finite level let \(n\ge4D_E+4\) be even.
For distinct lattices let \(a(L,L')\) denote their relative overlap. Then
\[
 |a(L,L')|\le [L':L\cap L']^{-n/2}.
\] The resulting matrix \(C\) has
a bounded inverse on bounded automorphic holomorphic frame functions
of the prescribed weight.
\end{lemma}

\begin{proof}
Let \(\mathfrak m\) be the level ideal, coprime to both
\(2\) and \(\bar{\mathfrak m}\), and write \(L_*\) for the reference
target lattice. A target lattice is commensurable with
\(L_*\) by multiplication in both directions by some
\(s\in\mathcal O_E\) prime to \(\mathfrak m\bar{\mathfrak m}\).
The shifted lattice is
\[
 L^{\sharp}=(\mathfrak m L:\bar{\mathfrak m}).
\]

For a source lattice \(A\), let
\(B_L=A\otimes L^{\sharp}\) be its self-dual tensor lattice
in the Weil space, where the tensor convention uses
conjugate source contractions. The normalized shifted
vacuum \(v_L\) is fixed by the operators
\[
 U(\delta)=\varepsilon(\delta)W(\delta)
\]
for
\(\delta\in B_L\). Here \(W\) is the Heisenberg action,
and \(\varepsilon\) is the common lattice semicharacter.
The level sum for \(L\) uses translates \(W(y)v_L\),
with \(y\) labeled by source rows modulo \(\mathfrak m L_*\)
after localization at \(\mathfrak m\).

The allowed labels
are the level subgroup in the first \(r\) rows and zero
in the remaining rows. The reduced matrices and their
unitary character give the same weights for all \(L\).
The two sums admit common representatives: the intermediate lattices agree at the level because the intervening scalars are prime to \(\mathfrak m\). The Chinese remainder theorem supplies representatives in their intersection. Denominator-lattice shifts fix the translate. Both the lift \(y\) and the difference \(d\) contain the factor \(2\) in the level quotient. Hence the vacuum semicharacter and their integral pairing phase equal one. Consequently
\[
 \varepsilon(d)=1,\qquad W(y+d)v_L=W(y)v_L.
\]


\proofstep{Orthogonality from the Heisenberg characters.}
Suppose \(\langle W(y)v_L,W(y')v_{L'}\rangle\ne0\). For every
\(\delta\in B_L\cap B_{L'}\), the two vectors are eigenvectors
of the same unitary \(U(\delta)\), with eigenvalue ratio
\(\exp(2\pi i\, E_A(\delta,y-y'))\) up to the convention for
the sign of the alternation \(E_A\). Put \(\Delta=y-y'\) and \(B_0=B_L\cap B_{L'}\).
Unitarity gives
\[
 \bigl(1-e^{2\pi i E_A(\delta,\Delta)}\bigr)
         \langle W(y)v_L,W(y')v_{L'}\rangle=0
       \qquad(\delta\in B_0).
\]
For a nonzero overlap this gives
\(\Delta\in B_0^\#\), with \(\#\) denoting the alternating
dual.

Choose \(s,s'\) prime to \(\mathfrak m\bar{\mathfrak m}\)
from the two comparisons with \(L_*\). The inclusion
\(ss' B_L\subseteq B_0\) implies
\[
 B_0^\#\subseteq(ss' B_L)^\#
             =(\overline{ss'})^{-1}B_L ,
\]
using \(B_L^\#=B_L\) and the conjugate adjoint rule. The source lattices in the
level sum contain \(N\mathcal O_E^n\) for an integer \(N\)
prime to a level integer in \(\mathfrak m\).
Contracting with \(Ne_k\) consequently gives
\(N\overline{ss'}(y_k-y'_k)\in L^{\sharp}\) in each row.
Take \(a\in\mathfrak m\), \(b\in\bar{\mathfrak m}\) with
\(a+b=1\). The definition of \(L^\sharp\) gives
\[
  bN\overline{ss'}(y_k-y'_k)\in\mathfrak m L,\qquad
  sbN\overline{ss'}(y_k-y'_k)\in\mathfrak m L_* .
\]
The scalar in the second inclusion is a unit modulo
\(\mathfrak m\), hence \(y_k,y_k'\) have the same level
label. This holds for every row.

For equal labels the same \(W(y)\) occurs in each argument.
Since it is unitary and the level-character weights have
modulus one, the normalized level contribution is
\(\langle v_L,v_{L'}\rangle\) times the common number of
labels. For its absolute value put
\(D'=B_{L'}/(B_L\cap B_{L'})\). The projection on the
\(B_{L'}\)-vacuum line, evaluated on \(v_L\), is
\[
      \Pi_{L'}v_L=\frac1{|D'|}\sum_{\delta\in D'}U(\delta)v_L .
\]
The terms are well defined because the denominator fixes
\(v_L\).

Distinct terms have different \(B_L\)-characters
by self-duality of \(B_L\), hence are orthogonal and each
has norm one. Locally the source lattice is free of rank
\(n\), so after the common shift the relative tensor index
is \(|D'|=[L':L\cap L']^n=m^n\). Therefore
\[
 |\langle v_L,v_{L'}\rangle|^2
   =\|\Pi_{L'}v_L\|^2=|D'|^{-1}=m^{-n}.
\]
The relative bound is the same for every source class. Summation
with the identical nonnegative stabilizer weights gives it for
the matrix \(C\).

\proofstep{Counting neighboring self-dual lattices.}
The number of self-dual \(L'\)
with index \(m\) is at most \(m^{2D_E}\). Set
\(J=L\cap L'\). Self-duality gives equal covolumes
and hence both indices over \(J\) equal to \(m\).
The intersection is an index-\(m\) sublattice of
\(L\). The number of index-\(m\) sublattices of
a free abelian group of rank \(D_E\) is at most
\(m^{D_E}\), by triangular Hermite normal form.
Indeed the recurrence in rank \(d\) is
\[
 a_d(m)=\sum_{k\mid m}k^{d-1}a_{d-1}(m/k)\le\sum_{k\mid m}k^{d-1}(m/k)^{d-1}=m^{d-1}\tau(m).
\]
There are at most \(m\) positive divisors, giving the asserted bound.

Now \(L'\) is an index-\(m\) sublattice of the
dual of \(J\), by the same covolume calculation.
The same count bounds the choices for \(L'\) for
each \(J\). Index one occurs only for \(L'=L\).
The normalized absolute row sum off the diagonal
is therefore at most
\[
       \sum_{m\ge2} m^{2D_E-n/2}
          \le \sum_{m\ge2} m^{-2}<0.65 .
\]
The estimate applies to every target lattice in the
system and every rank meeting the size condition.

\proofstep{The inverse and its multiplier.}
Let \(D\) be the positive diagonal of \(C\).
The unweighted bound on the off-diagonal row sum
makes \(D^{-1}C=1+T\) invertible on the bounded
automorphic functions by the series
\(\sum_{j\ge0}(-T)^j\). Each diagonal entry is
the fixed test-vector norm with the finite class
weights, so its inverse is bounded. The coefficient
matrix is compatible with the unit-modulus
automorphy multiplier on the target lattices.
It acts on bounded holomorphic frame functions
with the prescribed weight. The series converges
uniformly on compact subsets, since the weighted
supremum bound gives an ordinary supremum bound
on compact subsets of the frame domain. Thus its
inverse preserves holomorphicity and the multiplier.
Given the original bounded holomorphic family
\(F\), use \(F'=C^{-1}F\) in the unfolding.
\end{proof}

\paragraph{Reproduction on the thin slab.}

\begin{lemma}[Reproduction of the holomorphic frame functions]\label{lem:slab-reproducing-kernel}
For the Gaussian attached to the given stable-range weight, there is
a sufficiently thin slab and a constant \(\rho>0\) such that its
kernel reproduces each bounded holomorphic frame function of the
specified weight there. At an orthonormal frame \(x_0\) the scalar is
\[
 \rho=\frac{\int_{\rm slab}|K(x_0,y)|^2\,dy}{K(x_0,x_0)} .
\]
\end{lemma}

\begin{proof}
Use the Gaussian kernel with the row highest
polynomials in the prescribed weight. The stable
range contains these polynomials in the harmonic
tensor model. Take a finite family orthonormal
for the Fock pairing on this component. At a real
place use coordinates in which the Hermitian Gram
matrix of an orthonormal frame is
\(J=\operatorname{diag}(I_p,-I_q)\), \(p+q=r\).
The frame domain has invertible frames and
positive first \(p\) columns for this form. A thin
slab consists of frames with Gram matrix in a small
open neighborhood of \(J\), at every real place.
Indefinite Gram--Schmidt writes a frame as \(u b\)
in orthonormal coordinates, with \(u\) unitary for
\(J\) and \(b\) upper triangular.

For the first \(j\) columns of the frame \(y\), let
\(G_j=y_{[j]}^*J y_{[j]}\) and put \(\det G_0=1\). Since
\[
 G_j=b_{[j]}^*J_{[j]}b_{[j]},
\]
where the brackets on the square
matrices denote their leading blocks,
\[
 |b_{jj}|^2=\left|\frac{\det G_j}{\det G_{j-1}}\right|.
\]
On a sufficiently thin slab each \(\det G_j\) remains close to
\(\det J_{[j]}=\pm1\). The diagonal entries and their inverses
are therefore bounded. The Gram--Schmidt recursion then bounds
the remaining entries of \(b\) using these pivots and the
bounded entries of the Gram matrix.
The left unitary action preserves the slab and
Lebesgue measure because each unitary matrix has
determinant of absolute value one.

Let \(K(x,y)\) be the polarized Gaussian after
applying the highest polynomials in both arguments.
It is holomorphic in \(x\), and its right upper
triangular covariance is the weight character
\(\chi_\lambda\). For a unitary change \(u\), the two invariances are
\[
 K(y,x)=\overline{K(x,y)},\qquad K(ux,uy)=K(x,y).
\]
In the Gaussian
these changes preserve the Gram factors. Swapping
the two frames conjugates the polarized factors.
In the first argument the upper triangular change
contributes the determinant power of the Gaussian
and the diagonal powers of the two blocks of the
highest polynomial. Terms from same-argument
quadratic factors are killed by the harmonic
contractions.

At an orthonormal frame \(x_0\) the
diagonal Gaussian, in the harmonic letters \(z_l,w_l\),
is a positive determinant factor times
\(\exp(\sum_l t_l z_l w_l)\), with \(t_l>0\) from the
positive imaginary part of the polarized parameter.
For a nonzero polynomial \(P(z)=\sum_\alpha P_\alpha z^\alpha\)
in the chosen harmonic basis, coefficient extraction gives
\begin{equation}\label{eq:slab-gaussian-diagonal}
 \left.P(\partial_z)\overline P(\partial_w)
       \exp\!\left(\sum_l t_l z_l w_l\right)\right|_{z=w=0}
      =\sum_\alpha |P_\alpha|^2\alpha!\prod_l t_l^{\alpha_l}>0 .
\end{equation}
Only matching multi-indices survive the two differentiations.
Summing \eqref{eq:slab-gaussian-diagonal} over the basis and restoring the positive determinant
factor gives \(K(x_0,x_0)>0\).

\proofstep{The two torus characters.}
Consider the difference \(F\) of two functions with the same
value at \(x_0\), the prescribed right upper triangular weight
and the matching left diagonal unitary character. In coordinates
at \(x_0\), Gaussian elimination reduces a frame in the dense
cell to a lower triangular matrix \(L\). Define
\[
 D(z)=\operatorname{diag}(z,z^2,\ldots,z^r),
 \qquad L(z)=D(z)LD(z)^{-1}.
\]
The \((i,j)\) entry on or below the diagonal is \(z^{i-j}L_{ij}\),
so \(L(z)\) extends holomorphically to an invertible diagonal
matrix at zero.

For \(0<|z|\le1\), the first \(p\) diagonal entries
of \(D(z)\) have squared modulus at least \(|z|^{2p}\);
the others have squared modulus at most \(|z|^{2p}\).
Hence \(D(z)^*JD(z)\succeq |z|^{2p}J\), and the positive
Gram block stays positive after left multiplication. The
right diagonal changes that block by a congruence, so the
arc stays in the domain. For \(|z|=1\) the left and right
characters cancel and
\[
 F(L(z))=F(L).
\]
The maximum modulus
principle applied to \(F(L(z))-F(L)\) gives the equality at
zero. The value there vanishes by the prescribed weight and
\(F(x_0)=0\). The dense cell and continuity give \(F=0\).

For a fixed \(x_0\) this uniqueness expresses \(K(y,x_0)\)
as \(K(x_0,x_0)\) times the product of leading minors
\(\Delta_i\) of the coordinate matrices of \(y\) with
powers \(\lambda_i-\lambda_{i+1}\), with the weight
extended by zero at \(r+1\). This product has precisely
the required two characters. Possible negative powers
occur only at \(p,r\), where the minors do not vanish.
At a place with \(p,q>0\) the weights decrease within
each block and
\[
 \lambda_{p+1}-\lambda_p\ge n.
\]

On an orthonormal coordinate matrix \(U\), all the
leading minors are bounded by \(|\Delta_p(U)|\ge1\)
and \(|\det U|=1\). Write \(P=\mathrm{diag}(I_p,0)\) and \(Q=I-P\).
For \(U^*JU=J\) set \(N=PU+Q\) and \(N'=P+QU\).
Then
\[
 N^*N=U^*PU+Q=U^*(P-Q)U+U^*QU+Q=P+U^*QU=N'^*N'.
\]
The first
\(p\) columns of \(U\), written \(\binom{A}{C}\), satisfy
\(A^*A=I_p+C^*C\), so \(|\det N|=|\det A|\ge1\).
Writing \(W=N'N^{-1}\), we obtain
\[
 W^*W=(N^{-1})^*N'^*N'N^{-1}=(N^{-1})^*N^*NN^{-1}=I .
\] For \(N_i\) retaining the first
\(i\) rows of \(U\) and taking the other rows from the identity,
each row of \(N_iN^{-1}\) is a row of \(I\) or \(N'N^{-1}\).
All have norm one, so Hadamard gives
\[
      |\Delta_i(U)|=|\det N_i|\le|\det N|=|\Delta_p(U)|.
\]

The sum of the minor exponents up to \(r-1\) telescopes
to \(\lambda_1-\lambda_r\le-n\).
It follows that \(|K(y,x_0)|\) is bounded by a
constant times the product of \(|\Delta_p(y)|^{-n}\)
on the slab. The triangular factors there change
this bound only by a bounded factor. At a definite
place the slab is bounded and separated from singular
frames. A left unitary change and right triangular
covariance give the same bounds with any fixed first
frame \(x\).

\proofstep{Integrability in the noncompact directions.}
For integrability at an indefinite place write the
first \(p\) columns of \(y\) as \(\binom{A}{C}\).
The matrix
\[
 G=A^*A-C^*C
\]
stays in a bounded
neighborhood of \(I_p\) in the positive cone.
At fixed first columns use a left unitary isometry
to move them to \(\binom{G^{1/2}}{0}\). On the last
\(q\) columns this change has Jacobian one. The
cross Gram block bounds their top block and the
last Gram block then bounds their lower block.
Their total volume is bounded independently of
the first columns. For a square complex matrix \(A\),
polar coordinates have measure proportional to
\(d(A^*A)\) times compact unitary measure. This follows
from the triangular decomposition with positive
diagonal entries \(r_{ii}\), where both the matrix
and Gram Jacobians have the factor
\(\prod_i r_{ii}^{2(p-i)+1}\).

Since
\[
 A^*A=C^*C+G,
\]
integration in \(A\) amounts
to integration in a bounded set of \(G\) and is bounded
by a constant times \(\det(I_p+C^*C)^{-n/2}\).
To integrate \(C\), remove one row \(v\) and put
\(H=I_p+C_{\rm rem}^*C_{\rm rem}\). The determinant lemma
and the substitution \(v=wH^{1/2}\) give
\[
 \int_{\mathbb C^p}\det(H+v^*v)^{-a}\,dv
    =\det(H)^{1-a}
          \int_{\mathbb C^p}(1+\|w\|^2)^{-a}\,dw .
\]
The last integral is finite for \(a>p\). Starting with
\(a=n/2\), the last of the \(q\) removals has exponent
\(n/2-q+1>p\), since \(n/2\ge p+q\).
This proves absolute integrability of \(K(x,\cdot)\) on the
product slab. The pointwise bounds make the kernel bounded,
and the weighted frame-function norm bounds the ordinary norm
on the slab.


\proofstep{Evaluation by the torus average.}
Average a holomorphic frame function \(F\) against
\(\chi_\lambda^{-1}\) on the left diagonal unitary torus
at \(x_0\). Compactness gives a holomorphic average.
The average retains the right weight and has value
\(F(x_0)\) at \(x_0\). It also leaves the integral against
\(K(x_0,y)\) unchanged, by simultaneous invariance
and covariance in the first frame. Both this average
and \(K(y,x_0)\) have the two characters from the
uniqueness argument. Therefore
\[
 \int_{\rm slab}K(x_0,y)F(y)\,dy
 =\frac{\int_{\rm slab}|K(x_0,y)|^2\,dy}{K(x_0,x_0)}
   F(x_0).
\]
The numerator is finite and positive by boundedness,
integrability and continuity at \(x_0\). Move \(x_0\)
to any orthonormal frame by a left unitary change of
variables, and then to any frame by right weight
covariance. This proves reproduction with the same
positive constant \(\rho\).
\end{proof}

\paragraph{Unfolding the reproducing integral onto the original rows.}

\begin{proposition}[Spanning by the original theta rows]\label{prop:span-original-theta}
For the level and stable range \(n\ge4D_E+4\) of
Lemma~\ref{lem:theta-overlap-invertible}, unfolding
the doubled Weil kernel spans the bounded holomorphic target of the
prescribed weight by the original theta functions. At the base tuple write \(V_{\rm good}\) for the rational good row and \(\Theta_{\rm good}\) for the projected definite-source lift. After the prescribed scalar identifications,
\[
 \operatorname{im}\Theta_{\rm good}=V_{\rm good}.
\]
\end{proposition}

\begin{proof}
In the doubled Weil calculation the sum of the product
kernel over rational changes of target frame is
the finite sum of theta products, with the first
factor read at \((L,x)\) and the second at the other
target lattice and frame. On the unfolded side
the term for \(L'\) is
\[
    C(L,L')\int_{\rm slab} K(x,y)F'(L',y)\,dy .
\]
To obtain the doubled identity, apply the definite
Hermitian Siegel--Weil identity to the admissible
doubled test function at zero in the auxiliary
variables. The difference from the Eisenstein sum
solves the doubled heat equation. The invariant
sum of highest-polynomial differential operators
therefore also kills that difference.

In the
Eisenstein sum a maximal isotropic subspace of
the doubled space is either a unitary graph or
contains a nonzero isotropic vector in one summand.
Writing \(I(U)\) for the term attached to a maximal isotropic space \(U\), the no-coinvariants condition and the support of the finite Weil pairing give, respectively,
\[
 I(U)=0\quad\text{if }U\text{ is not a graph},
 \qquad I(\operatorname{graph}g)=0\quad\text{if }g\text{ is outside the level group}.
\]
On a graph the finite pairing is the overlap of the two test functions. For the
remaining graphs the Weil action on the translated
frame and the conjugate multiplier give the
integrand above. The double sum and
integration may be interchanged.

For fixed \(x\),
the Gaussian at \(x\) is integrable in \(y\) on a
sufficiently thin slab, the functions \(F'\) are
uniformly bounded on that slab, and
\[
 \sum_{L'}|C(L,L')|<\infty
\]
by the overlap
estimate. The group changes preserve the measure
and the slab, and the unit-modulus multiplier
cancels against the conjugate in the theta product.
Integrate over the fundamental domain for these
group changes. The finite sum of theta products
then yields a finite sum of theta functions in
\((L,x)\), whose scalar coefficients are the integrals
of the second factors against \(F'\).
The unfolded expression is
\[
    \rho\sum_{L'} C(L,L')F'(L',x)=\rho F(L,x).
\]
Dividing the coefficients by \(\rho\) gives the
theta spanning assertion. Only the strict bound
on \(T\) is used for an estimate on the rank.
The overlapping lattice system and the doubled
kernel both use the prescribed finite level in this
application. \label{cor:theta-row-rank-1}

\proofstep{Identification of the lowest-vector realization.}
In the adelic realization, evaluate the target
discrete-series map on its nonzero lowest weight
vector. If \(v_0\) is that lowest vector, cyclicity gives the test for an intertwiner \(A\):
\[
 Av_0=0\quad\Longrightarrow\quad A(U(\mathfrak g)v_0)=0\quad\Longrightarrow\quad A=0.
\]
Here \(\mathfrak g\) denotes the complexified Lie algebra on the Harish--Chandra module.
The bottom evaluation has the given compact
torus weight and is killed by the lowering
operators. After extracting the common
determinant character it gives the holomorphic
frame function of the preceding argument. Cusp
rapid decay gives its weighted boundedness. The
nonzero finite Weil vector of the base pair is
a finite translate of the level test vector for
that argument. Conjugating the frames and
undoing this translation preserve the allowed
harmonic source span. The lowest-vector identity spans the unprojected range of \(\Theta\): unfold each term as a definite-source lift summed over source classes. For the good-type projector \(e\),
\[
 \operatorname{im}(e\Theta)=e(\operatorname{im}\Theta).
\]

\proofstep{Projection onto the good row.}
Spanning before taking finite types gives access to all row coordinates. The determinant at the selected place uses the good isotypic component, so the projection has to be applied to the same spanning family.

Now take the base tuple on the actual row. The
trivial auxiliary character and empty auxiliary
level identify its projected lift with the class
sum in the row construction. For its finite unipotent group \(U\) and character \(\chi\),
the projection is
\[
 e_\chi=|U|^{-1}\sum_{u\in U}\chi(u)^{-1}u.
\]

Changing variables \(v=u_0u\) in the sum yields
\[
 u_0e_\chi=|U|^{-1}\sum_{v\in U}\chi(u_0^{-1}v)^{-1}v
            =\chi(u_0)e_\chi .
\]
Conversely \(uf=\chi(u)f\) implies \(e_\chi f=f\). For the equivariant theta map \(T\),
\[
                  T(e_{\chi,\mathrm{src}}v)
                         = e_\chi T(v),
\]
where the source operator uses the dual transformation
rule of the row reading. This imposes that rule on each
preimage of a good row vector. The equivariance of the row reading
makes the average fix the given target vector.
On the source the average imposes the dual
unipotent rule. Thus a good row with this type
is the reading of a class sum of functions on
the source with the base finite transformation
law.

On a finite list of source classes, the
harmonic coefficient lattice spans its allowed
value space over \(\mathbb C\). For the class
stabilizer, the finite trace with its level
multipliers preserves that lattice. For a stabilizer \(H\) on the complex value space, its trace and projector satisfy
\[
 \operatorname{Tr}_H=\sum_{h\in H}h,\qquad |H|^{-1}\operatorname{Tr}_H=\operatorname{pr}_{H\text{-invariants}}.
\]
The operators include the specified level multipliers. Using the unnormalized trace on
integral vectors therefore spans all such class
functions. Their class sums span the given good
row.

A summand at one source class may be moved
to a finite split-supported representative
away from the excluded places. Rational
automorphy and the finite transformation
rule transport its polynomial in the harmonic
lattice. The target character factor is
absorbed by that rule. Each resulting term
is the complex reading of an admitted original
theta datum. This uses the split-support and
level conditions in the definition of the
allowed original data, and the base-series
rational comparison on the same row.

Finally, any vector of the rational good row
whose complex reading is in the span of those
data lies in their rational span. Indeed if \(W\) is their rational span, extension of
scalars over a field is exact, so the kernel of
\[
         V_0\longrightarrow (V_0/W)\otimes_K\mathbb C
\]
is \(W\). The comparison identifies the complex reading
with this scalar extension. Applying the type projection, which
fixes the good row, gives the full span of the
projected original theta elements in \(V_0\).
The finite set of source classes and the
finite harmonic lattices give a finitely generated
row module.

Equivalently, in the rational
comparison the original theta elements lie
in the fixed finite level lattice. Over its DVR,
a generating set for a finite module contains
a subset projecting to a basis modulo the
maximal ideal. Since this module is torsion-free,
the chosen generators form a lattice basis.
This proves the original full-rank comparison
on the fixed tied data and in the stated stable
range. Its use for the division rank later
also requires the identification with the
selected point's period and row dimension.
\label{cor:theta-row-rank-2}
\end{proof}

\paragraph{Containment for the twisted rows.}
Fix the tested character at one admitted auxiliary
level, and work over its coefficient DVR \(O_k\).
The finite type multiplier on the source level
group is a product of its base character, the
splitting character at the two type primes, and
the local components of the fixed and varying
twists. The base character agrees with the
given row character there. All its values are
roots of unity in the coefficient field.
At a place over a type prime \(g_i\) the field
assumption includes the roots of order dividing
\(g_i(g_i-1)\), since the host places there have
residue field \(\mathbb F_{g_i}\).
At the distinguished and auxiliary places use
the prescribed character values in \(O_k\).
If \(\mu\) denotes the product character and \(u\) a level
element, each of its factors has finite order in the admitted
coefficient field. Thus the normalization is
\[
       \mu(u)=\mu_{\rm base}(u)\mu_{\rm spl}(u)
                          \mu_{\rm fixed}(u)\mu_{\rm aux}(u)
                       \in O_k^\times .
\]

For a twisted generator the source translate
differs from its original representative by a
change supported at the chosen split places.
Those places avoid the coefficient prime and
the fixed exclusions. In the harmonic coefficient
lattice the two readings agree over \(O_k\) after
this change. At places over the coefficient prime
the self-dual lattices agree, and all rational
denominators of the other changes are units there.
Write \(h\) for the resulting finite source point
and \(P\) for its integral harmonic polynomial.
The finite trace of \(P\) on the double coset
of \(h\) is the sum of its rational source
translates with the finite level multipliers.

At a given source point the summands are indexed
by a coset of a finite stabilizer, so this is
a finite sum in its harmonic coefficient lattice.
The reindexing uses the bijection of finite indexing sets
\[
 H\longrightarrow h_0H,\qquad h\longmapsto h_0h,
\]
where \(H\) is the stabilizer and \(h_0H\) is the coset indexing a translated sum. Including the level multiplier gives the dual transformation rule. It belongs to the integral module
of class functions for this level.

Unfold its class sum using the finite source
representatives and the same stabilizer weights
as in the original theta formula. On the kernel,
rational left translation fixes the value at
zero and the right level action is the inverse
of the multiplier in the trace. The two factors
cancel.

For one contributing class let \(n_h\) be its stabilizer
order and \(\Theta_h(P)\) the value of the corresponding
flat series term. Each of the \(n_h\) translated terms
has that same value after the two finite multipliers cancel.
The resulting term in the class sum is
\[
       n_h^{-1}\sum_{i=1}^{n_h}\Theta_h(P)=\Theta_h(P).
\]
The reciprocal arises only in the characteristic-zero
class-sum evaluation, after forming the integral trace. The result is the desired single flat
twisted series, multiplied by its norm-one
determinant character on the source point.
That scalar is a root of unity in \(O_k\).
Including it in the integral trace therefore gives
the correct twisted generator. The trace takes
no inverse stabilizer order on the integral
module.

By the comparisons on the fixed row, the class
sum is the reading of an element of the finite
level lattice. The twisted generator has the same
reading. Injectivity over the extension field
identifies these vectors, so the generator belongs
to that lattice before applying the good-type
projection. This proves the twisted containment
used in the signed-index comparison. It applies
to both sides and every admitted character of
the specified depth. \label{cor:theta-row-rank-3}
